\section{Discusión}
\label{sec:discusion}
% Interpretación, comparación con la teoría, hallazgos inesperados, limitaciones. Acá se
% responden, sin numerarlas, todas las preguntas sueltas de la consigna (docs/REQUISITOS.md §1):
%   - efecto de disminuir h; compromiso precisión-costo; si un paso más chico siempre conviene;
%   - cuál método tiene menor error, dependencia de h, operaciones por iteración, si el más
%     preciso siempre conviene;
%   - qué representa k, cuál organización remedia más rápido, cómo cambia la gráfica con k, si k
%     alto es mayor capacidad de respuesta, qué factores reales lo explican;
%   - hk > 1, valores negativos, sentido de una cantidad negativa de equipos;
%   - λ: disminuir, crecer o estabilizarse; interpretación para una organización real;
%   - T1190 y la superficie de ataque; relación entre k y el proceso de remediación;
%   - limitaciones del modelo (las siete preguntas de la sección 12 de la consigna);
%   - análisis global: precisión, tamaño del paso, costo, estabilidad, interpretación.

[PENDIENTE]
